\documentclass[11pt,a4paper]{article}
\PassOptionsToPackage{pdfusetitle=false}{hyperref}
\usepackage[utf8]{inputenc}
\usepackage[T1]{fontenc}
\usepackage[margin=2.5cm]{geometry}
\usepackage{amsmath,amssymb}
\usepackage{booktabs,array}
\usepackage{xcolor}
\usepackage{tcolorbox}
\tcbuselibrary{skins,breakable}
\usepackage{enumitem}
\usepackage[pdfusetitle=false]{hyperref}
\usepackage[numbers]{natbib}
\bibliographystyle{unsrtnat}
\hypersetup{colorlinks=true,linkcolor=blue!60!black,urlcolor=blue!60!black,
            citecolor=blue!60!black,
            pdftitle={New Physics Predictions},
            pdfauthor={Nova Spivack}}
\setlength{\emergencystretch}{2em}
\usepackage{tikz}
\usetikzlibrary{arrows.meta,positioning,shapes.geometric,calc,fit,backgrounds}

%─── Color palette ───────────────────────────────────────────────────────────
\definecolor{boxblue}{RGB}{220,235,255}
\definecolor{boxorange}{RGB}{255,240,210}
\definecolor{boxgreen}{RGB}{220,245,220}
\definecolor{boxred}{RGB}{255,230,225}
\definecolor{boxpurple}{RGB}{240,225,255}
\definecolor{boxyellow}{RGB}{255,252,210}

%─── Macros ──────────────────────────────────────────────────────────────────
\newcommand{\term}[1]{\textbf{#1}}
\newcommand{\lean}[1]{%
  \penalty-100
  \begingroup\def\_{\textunderscore\allowbreak}\texttt{#1}\endgroup}
\newcommand{\CatAL}{\textbf{CatAL}}
\newcommand{\CatAD}{\textbf{CatAD}}
\newcommand{\CatA}{\textbf{CatA}}
\newcommand{\PHIMDL}{$\Phi_{\rm MDL}$}
\newcommand{\PhiMDL}{\Phi_{\rm MDL}}

\newcommand{\TutorialDOI}{10.5281/zenodo.20669152}
\date{2026}

\title{\textbf{New Physics Predictions}\\[6pt]
\large What GTE Predicts Will and Won't Be Found\\[8pt]
\normalsize
\href{https://novaspivack.com/research}{novaspivack.com/research}\texorpdfstring{\quad}{ }
\href{https://doi.org/10.5281/zenodo.20560550}{P48 complete theory monograph}\\[4pt]
\small\href{https://doi.org/\TutorialDOI}{doi.org/\TutorialDOI}
  \textit{(registration pending)}\\[4pt]
\normalsize\textit{UGP Physics Tutorial Series}}
\author{Nova Spivack}

\begin{document}
\setcounter{tocdepth}{2}
\maketitle
\tableofcontents
\newpage

\begin{tcolorbox}[colback=blue!6,colframe=blue!40!black,
  title=\textbf{UGP Physics Tutorial Series}]
\textbf{Prerequisites (read first):}
  \href{https://doi.org/10.5281/zenodo.20657898}%
    {\emph{How the Standard Model Quantum Numbers Emerge}}
  $\cdot$
  \href{https://doi.org/10.5281/zenodo.20669144}%
    {\emph{Forces, Interactions, and the V$-$A Structure}}
  $\cdot$
  \href{https://doi.org/10.5281/zenodo.20657885}%
    {\emph{The $\Phi_{\rm MDL}$ Field: From Discrete Certificate to Continuum Physics}}\\[4pt]
\textbf{Next tutorials:}
  \href{https://doi.org/10.5281/zenodo.20669142}%
    {\emph{Falsifiability: How to Test the GTE Framework}}\\[4pt]
\textbf{See also:}
  \href{https://doi.org/10.5281/zenodo.20657889}%
    {\emph{The GTE Polynomial: A Step-by-Step Cheat Sheet}}
  $\cdot$
  \href{https://doi.org/10.5281/zenodo.20657872}%
    {\emph{The Levels of the Theory}}
\end{tcolorbox}

\begin{tcolorbox}[colback=boxgreen,colframe=green!50!black,
  title=\textbf{Claim-Strength Taxonomy}]
Throughout this tutorial, results are labeled with a confidence grade:
\begin{itemize}[itemsep=2pt]
  \item \textbf{$\mathbf{A}_{\rm Lean}$ (CatAL)} --- Machine-certified in Lean~4 with zero \texttt{sorry} and zero custom axioms.  Independently verifiable by anyone with Lean installed.
  \item \textbf{A/D (CatAD)} --- Analytically derived with full derivation, computationally confirmed.  Not yet machine-certified.
  \item \textbf{A (CatA)} --- Analytically derived; derivation complete but not independently machine-certified.
  \item \textbf{A/D (conditional)} --- Derived under a named, stated premise; the premise is itself an open problem or a CatB assumption.
  \item \textbf{B (CatB)} --- A stated working hypothesis or structural premise; not yet derived.
\end{itemize}
\end{tcolorbox}

%══════════════════════════════════════════════════════════════════════════════
\section{Why New Physics Predictions Matter}
\label{sec:intro}
%══════════════════════════════════════════════════════════════════════════════

\begin{tcolorbox}[colback=boxblue,colframe=blue!40!black,
  title=The core argument]
With zero free parameters, every numerical output of the GTE framework is a
\emph{genuine prediction} that the framework could have failed but did not.
The same zero-parameter structure that derives the Standard Model's known
physics also constrains what lies beyond it.  This tutorial catalogs what
GTE forbids, what it predicts, and why those predictions are accessible to
current and near-future experiments.
\end{tcolorbox}

Since the Large Hadron Collider began operation, it has confirmed the Higgs
boson but has found no evidence of physics beyond the Standard Model (BSM):
no supersymmetric particles, no extra dimensions, no new forces.  Many BSM
theories predicted new particles that should already have appeared.  They have not.

\term{GTE} (\term{Generative Triple Evolution}) is designed differently.  Rather
than adding new ingredients to the Standard Model, it derives the SM from a
single 19-bit algebraic certificate: the polynomial
$p(L,C,R) = C + R - CR - LCR$ over $\mathrm{GF}(7)$.  This certificate has no
free parameters.  Its zero-parameter architecture means that any structure not
forced by the certificate is absent — and that is precisely why GTE predicts
the LHC has found nothing.  The theory explains the null result.

But GTE also makes specific predictions about what \emph{will} be found.  These
predictions are sharp: each has a named experiment, a numerical threshold, and
an explicit falsification criterion.  This tutorial presents each one with its
physical mechanism, its numerical value, and what observation would confirm or
refute it.

%══════════════════════════════════════════════════════════════════════════════
\section{What GTE Forbids: The Structural Absences}
\label{sec:forbidden}
%══════════════════════════════════════════════════════════════════════════════

\subsection{No Supersymmetry}

\term{Supersymmetry} (SUSY) introduces partner particles (squarks, gluinos,
neutralinos) for every SM particle.  In GTE, the complete particle spectrum is
determined by the $\mathbb{Z}_7$ winding sectors of the $\Phi_{\rm MDL}$ field.
There are exactly five PSC-admissible winding sectors in the SM branch
($w \in \{0,2,3,4,6\}$), plus the dark-mirror branch ($w \in \{1,5\}$).  No
additional sectors are available.  Supersymmetric partners would require a
doubling of the winding spectrum with no mechanism to generate it in the
GTE certificate structure.

GTE does not forbid SUSY by a theorem that explicitly names SUSY; rather, the
complete particle ontology is derived from $\mathbb{Z}_7$ arithmetic and admits
no room for superpartners \cite{SpivackGTECompleteFramework}.  Every null SUSY
result from the LHC is consistent with this prediction.

\subsection{No Axion}

The \term{axion} is a hypothetical particle proposed by Peccei and Quinn to
explain why the strong CP angle $\theta_{\rm QCD}$ is observed to be less than
$10^{-10}$ despite the SM having no reason to favor $\theta = 0$.

In GTE, the strong CP problem is \emph{structurally resolved} without any axion.
The internal symmetry group $F_{21} = \mathbb{Z}_7 \rtimes \mathbb{Z}_3$ has
abelianisation $F_{21}^{ab} = \mathbb{Z}_3$ — pure color, zero CP-odd phase.
This group-theoretic fact forces $\theta_{\rm QCD} = 0$ exactly, with three
independent proofs, all machine-certified (\CatAL)
\cite{SpivackGTECompleteFramework}.

\begin{tcolorbox}[colback=boxred,colframe=red!60!black,
  title={Strong CP Prediction (\CatAL, three independent proofs)}]
$\theta_{\rm QCD} = 0$ exactly from $F_{21}$ group theory.  No axion is
predicted and no axion mechanism is needed.  Every null axion search (ADMX,
CASPEr, BabyIAXO, HAYSTAC) is therefore a \textbf{positive confirmation} of
the GTE structure, not a failure to find something.
\end{tcolorbox}

\textbf{Falsification:} A definitive axion detection at any mass and any
coupling strength would falsify the GTE strong-CP mechanism.

\subsection{No Large Extra Dimensions}

Theories with large extra dimensions (e.g., Arkani-Hamed, Dimopoulos, Dvali)
explain the hierarchy problem by diluting gravity into higher-dimensional space.

In GTE, the $3{+}1$-dimensional spacetime structure is not postulated — it is
derived.  The CMCA (\term{Chiral Minkowski Cellular Automaton}) has exactly one
non-trivial persistent-glider sector, which lives in $1{+}1$ dimensions.
The continuum $\Phi_{\rm MDL}$ field then lifts this to $3{+}1$ dimensions via
the Algebraic Lifting Theorem (\CatAL).  A $2{+}1$-dimensional persistent-glider
sector analogous to the SM is absent by theorem (the ``dimensional staircase''
result: no canonical $2{+}1$D discrete substrate compatible with the SM orbit
exists; \CatAD) \cite{SpivackGTECompleteFramework}.  No room is available for
compact or large extra dimensions within the framework.

\subsection[No New Light Particles Below Lambda-GTE]{No New Light Particles Below $\Lambda_{\rm GTE}$}

The GTE certificate structure defines a natural energy scale:
\begin{equation}
  \Lambda_{\rm GTE} = 7 \times M_{\rm kink} = \frac{8}{7}\,m_\tau
  \approx 2.03~\text{GeV}.
  \label{eq:lambda}
\end{equation}
Below $\Lambda_{\rm GTE}$, the EFT (\term{effective field theory}) description
is complete: the $\mathbb{Z}_7$ winding sectors are all accounted for, and no
new stable structures can appear in the SM branch.  Any new light stable particle
below $\Lambda_{\rm GTE}$ that is not in the SM spectrum would require a new
winding sector, which the certificate structure does not permit.

This does \emph{not} mean the EFT is trivial above $\Lambda_{\rm GTE}$.  The
kink dissolution threshold at $\Lambda_{\rm GTE}$ is a new-physics scale with
specific signatures (see \S\ref{sec:lambda_collider}).

%══════════════════════════════════════════════════════════════════════════════
\section{What GTE Predicts: The Affirmative Case}
\label{sec:predictions}
%══════════════════════════════════════════════════════════════════════════════

\subsection{The Mirror Branch Dark Sector}

The $\mathbb{Z}_7$ arithmetic that generates the SM winding sectors
$\{0,2,3,4,6\}$ does not exhaust $\mathbb{Z}_7$.  The complementary winding
sectors $\{1,5\}$ are as consistent with $\mathbb{Z}_7$ arithmetic as any SM
sector.  They form the \term{mirror branch} — a dark sector forced by the same
arithmetic that generates the SM \cite{SpivackGTECompleteFramework}.

\begin{tcolorbox}[colback=boxorange,colframe=orange!60!black,
  title={Definition: Mirror Branch (\CatAD)}]
The \emph{mirror branch} consists of $\mathbb{Z}_7$-admissible three-tape
configurations with primary winding number $w = 5$ (the $\mathbb{Z}_7$
conjugate of the SM lepton sector, $w = 2$).  Mirror particles carry
SU(3)$_{\rm dark}$ colour and \textbf{no SM gauge charges}: no electric charge,
no SM weak charge, no SM color.  They are dark by construction — the winding
algebra forbids SM-gauge coupling at dimension-4.
\end{tcolorbox}

The mirror branch has three generations, for the same combinatorial reason that
the SM has three: the PSC enumeration that forces $N_{\rm gen} = 3$ applies to
any $\mathbb{Z}_7$-admissible sector.

\subsection{The Mirror-Branch Mass Spectrum}

The mass spectrum of the mirror branch is computed by applying the same cascade
formula that gives SM lepton masses, with mirror-branch winding inputs.  The
three mirror lepton masses are \cite{SpivackGTECompleteFramework}:
\begin{align}
  m_{\chi_1} &= 0.54~\text{MeV} \qquad \text{(Generation 1, primary dark matter candidate)},\\
  m_{\chi_2} &= 24.5~\text{MeV} \qquad \text{(Generation 2)},\\
  m_{\chi_3} &= 3.60~\text{GeV} \qquad \text{(Generation 3)}.
\end{align}
These values are \CatAD.  The generation-2/1 mass ratio is
$m_{\chi_2}/m_{\chi_1} = 45.4$ — a prediction: if a dark lepton signal is found
at any one of these masses, the theory predicts where to find the other two.

The dark sector has a gauge group SU(3)$_{\rm dark}$, forced by the same
Group of Elegance mechanism that selects QCD.  There is \emph{no dark photon}:
the absence of a U(1) factor in the Elegant Kernel for the mirror winding sector
is Lean-certified.

\begin{tcolorbox}[colback=boxred,colframe=red!60!black,
  title={No Dark Photon (Lean-certified)}]
GTE predicts no dark photon at any mass.  Every null result from dark-photon
searches (NA64, BaBar, Belle~II kinematic scan, NA62) is a
\textbf{positive confirmation} of the mirror-branch structure.  Any dark photon
discovery at any mass would falsify the GTE dark sector model.
\end{tcolorbox}

\subsection[The Dark Matter Candidate: chi1 at 0.54 MeV]{The Dark Matter Candidate: $\chi_1$ at 0.54 MeV}

The lightest mirror lepton $\chi_1$ at 0.54~MeV is the primary dark matter
candidate.  Its spin-independent cross-section with SM nucleons is predicted
to be $\sigma_{\rm SI} \sim 10^{-61}$~cm$^2$ (\CatAD)
\cite{SpivackGTECompleteFramework} — approximately 30 orders of magnitude below
any current or projected direct-detection sensitivity.  This is not a tuned
value; it follows from the winding isolation at dimension-4.

This dark matter is \term{ultra-faint interacting} (UFIDM): it interacts
gravitationally and via SU(3)$_{\rm dark}$, but is invisible to WIMP-style
searches.  The predicted relic density is
$\Omega_{\rm DM} h^2 = 0.11994$ (\CatAD), matching the Planck~2018
measurement to $-0.15\%$ without any free parameters.

%══════════════════════════════════════════════════════════════════════════════
\section{GTE-P7: The Highest-Priority New-Physics Signal}
\label{sec:gtep7}
%══════════════════════════════════════════════════════════════════════════════

\begin{tcolorbox}[colback=boxyellow,colframe=orange!60!black,
  title={GTE-P7 Prediction (\CatAD) --- Belle II Target}]
A dark lepton resonance at $m_{\rm GTE\text{-}P7} = 211.9$~MeV is predicted
by the mirror-branch winding sector of the GTE polynomial
\cite{SpivackGTECompleteFramework}.  This is the highest-priority new-physics
prediction accessible to running experiments.
\end{tcolorbox}

\subsection{What GTE-P7 Is}

GTE-P7 is not $\chi_1$, $\chi_2$, or $\chi_3$, but rather a resonance
appearing in the Elegant Kernel computation for the mirror branch.  Its mass
arises from the same UCL (Universal Calibration Law) pipeline that generates the
SM lepton masses, applied to a specific mirror-branch cascade triple.

The numerical value $211.9$~MeV is \CatAD: the derivation chain is complete
through the mirror-branch winding sector.

\subsection{Where to Look}

The GTE-P7 mass lies in the range accessible to Belle~II
\cite{SpivackGTECompleteFramework}, currently taking data, in the channels:
\begin{itemize}[itemsep=3pt]
  \item $e^+e^- \to \chi_? + \text{invisible}$ — a single charged track with
    missing energy, characteristic of a dark lepton coupled via a portal.
  \item Displaced-vertex signatures — if the dark lepton decays with a
    centimeter-scale decay length inside the detector.
\end{itemize}
The full Belle~II dataset (50~ab$^{-1}$, projected by 2031) will cover the
signal window $[200, 225]$~MeV with sufficient sensitivity to confirm or exclude.

\textbf{Falsification:} The GTE-P7 prediction is falsified if no resonance is
observed in the 200--225~MeV window after the full 50~ab$^{-1}$ dataset, at
the expected coupling scale.

\subsection{Why This Mass and Not Another}

The cascade formula that places GTE-P7 at 211.9~MeV has no free parameters once
the mirror-branch winding sector is specified.  The same pipeline that correctly
derives $m_e$, $m_\mu$, $m_\tau$ at $0.295\%$ RMS is applied to the
mirror-branch triple.  The mass prediction is a corollary of the same arithmetic
that governs the SM masses.  A different mass would require either a different
winding sector or different cascade arithmetic — neither of which the framework
permits.

%══════════════════════════════════════════════════════════════════════════════
\section{Proton Stability: Dimension-4 Absolutely Forbidden}
\label{sec:proton}
%══════════════════════════════════════════════════════════════════════════════

\subsection{What the SM and GUTs Predict}

Grand Unified Theories predict proton decay via the channel
$p \to e^+ \pi^0$, a dimension-6 operator suppressed by $M_{\rm GUT}^{-2}$,
giving a proton lifetime $\tau_p \sim 10^{34}$--$10^{36}$~yr.
This is the primary target of Hyper-Kamiokande (under construction) and DUNE.

\subsection{The GTE Result: Dimension-4 Forbidden, Dimension-6 Not}

In GTE, baryon number is a \emph{topological charge} of the $\mathbb{Z}_7$
winding field:
\begin{equation}
  B = \frac{1}{3}\sum_j \chi_q(w_j),
\end{equation}
where the sum is over the three constituent kink windings.  Topological charges
are exactly conserved: they cannot change under any local deformation of the
field.

This makes dimension-4 baryon-number violation not merely suppressed but
\emph{structurally impossible} in GTE.  Machine-certified:
\lean{gte\_baryon\_number\_topological\_charge}
(\CatAL, zero sorry, \texttt{Universality/GUTStructure.lean})
\cite{SpivackGTECompleteFramework}.

\begin{tcolorbox}[colback=boxgreen,colframe=green!50!black,
  title={Proton Stability Theorem (\CatAL)}]
Dimension-4 proton decay is structurally forbidden: baryon number is a
topological winding invariant that cannot be violated by any local operator of
dimension $\leq 4$.  The proton is topologically stable against dimension-4
operators.\\[4pt]
\textbf{Note:} Dimension-6 operators (GUT-scale suppressed, $\sim M_{\rm GUT}^{-2}$)
are \emph{not} forbidden by this theorem.  The standard GUT-predicted channel
$p \to e^+ \pi^0$ may exist at the dimension-6 level.
\end{tcolorbox}

This is one of GTE's sharpest predictions.  A single verified proton decay
event identified as a dimension-4 process would falsify the framework.  Current
experimental limits already place the proton lifetime at
$\tau_p > 1.6 \times 10^{34}$~yr for the dominant $e^+\pi^0$ channel (Super-K).

\textbf{What the Lean proof establishes:}
The theorem \lean{gte\_baryon\_number\_topological\_charge} certifies that $B$
is conserved as a topological charge of the $\mathbb{Z}_7$ winding field.  This
means no perturbative operator, no SM radiative correction, and no GUT-scale
coupling can violate $B$ at dimension-4 in the GTE EFT.

%══════════════════════════════════════════════════════════════════════════════
\section[The Lambda-GTE Scale as a Collider Target]{The $\Lambda_{\rm GTE}$ Scale as a Collider Target}
\label{sec:lambda_collider}
%══════════════════════════════════════════════════════════════════════════════

The kink dissolution threshold
\begin{equation}
  \Lambda_{\rm GTE} = 7\,M_{\rm kink} = \frac{8}{7}\,m_\tau \approx 2.03~\text{GeV}
\end{equation}
is the energy at which the GTE EFT description starts to break down: above
this scale, individual kinks dissolve into the $\Phi_{\rm MDL}$ substrate and
new topological structure can form.

\begin{tcolorbox}[colback=boxblue,colframe=blue!40!black,
  title=What to look for at $\Lambda_{\rm GTE}$]
\begin{itemize}[itemsep=3pt]
  \item A change in the effective number of degrees of freedom in hadronic
    cross-sections near $\sqrt{s} \approx 2$~GeV.
  \item Anomalous energy deposition patterns characteristic of kink dissolution
    in the $\Phi_{\rm MDL}$ field.
  \item Lattice QCD calculations showing a transition in the kink structure
    function around $\Lambda_{\rm GTE}$.
\end{itemize}
\textbf{If no structure appears:} The absence of any new hadronic feature below
$\sim 3$~GeV would be informative but not immediately falsifying, since the
transition may be gradual.
\end{tcolorbox}

The $\Lambda_{\rm GTE}$ scale is first-principles derived, not a tunable cutoff.
It is fixed by $m_\tau$ and $|\mathbb{Z}_7| = 7$, both of which are themselves
derived.  This makes it a genuine prediction, not a parameter.

%══════════════════════════════════════════════════════════════════════════════
\section{Gravitational Wave Predictions}
\label{sec:gw}
%══════════════════════════════════════════════════════════════════════════════

\subsection[Primordial Gravitational Waves: r = 0 Exactly]{Primordial Gravitational Waves: $r = 0$ Exactly}

The GTE inflationary mechanism produces scalar perturbations with spectral tilt
$n_s = 1 - \ln 2/(2\pi^2) = 0.96488$ (\CatAL), confirmed at $-0.005\sigma$
by Planck~2018 \cite{SpivackGTECosmology}.

For the tensor-to-scalar ratio $r$, GTE predicts:
\begin{equation}
  r = 0 \qquad \text{(exactly, } \CatA\text{).}
\end{equation}
Primordial gravitational waves are absent because the inflationary mechanism in
GTE does not involve an inflaton field with large-field excursions — the perturbation
spectrum is generated by $\mathbb{Z}_7$ defect formation, not by a slow-roll
inflaton.

\textbf{Current sensitivity:} Planck+BICEP/Keck place $r < 0.036$ (95\% CL).
CMB-S4 (projected $\sim$2028) will reach $\sigma(r) \approx 0.002$.  A detection
of $r > 0.01$ at $>3\sigma$ would falsify the GTE inflationary mechanism.

\subsection[The Z7 Phase Transition Gravitational Wave Signal]{The $\mathbb{Z}_7$ Phase Transition Gravitational Wave Signal}

The $\mathbb{Z}_7$ symmetry breaking at temperature $T_G \approx 0.7$~GeV
(the domain-wall ordering crossover) produces a network of domain walls that
annihilates by $T_{\rm ann} \approx T_G$ (wall lifetime $\sim 5 \times 10^{-24}$~s).
This rapid annihilation generates a gravitational wave background peaked near
the frequency corresponding to $T_G$ \cite{SpivackGTECompleteFramework}.

\begin{tcolorbox}[colback=boxorange,colframe=orange!60!black,
  title=GW Background from $\mathbb{Z}_7$ Phase Transition (\CatAD)]
A gravitational wave background from the $\mathbb{Z}_7$ domain-wall network is
predicted in the frequency band accessible to the Einstein Telescope and LISA.
The zero-relic prediction (complete annihilation before BBN) is a structural
GTE result; any detection of a wall relic density above the $T_G$ level would
falsify the annihilation mechanism.
\end{tcolorbox}

%══════════════════════════════════════════════════════════════════════════════
\section{The Zero-Free-Parameter Principle: Every Output Is a Prediction}
\label{sec:zfp}
%══════════════════════════════════════════════════════════════════════════════

\begin{tcolorbox}[colback=boxpurple,colframe=blue!40!black,
  title=The key insight about falsifiability]
In a theory with free parameters, a measurement that disagrees with a prediction
can always be accommodated by retuning the parameters.  In GTE, there are no
free parameters to retune.  Every output — the nine charged-fermion masses, the
Weinberg angle, the dark matter mass, the proton lifetime, $r$, $n_s$ — is
derived, not fitted.  A single confirmed disagreement at $>3\sigma$ precision
would refute the theory as currently formulated, with no escape route.
\end{tcolorbox}

This makes GTE \emph{more} falsifiable than the Standard Model, not less.  The
SM has 19 free parameters (plus at least 3 for neutrinos) that can absorb any
discrepancy.  GTE has zero.  The absence of new LHC physics is not a negative
result for GTE — it is a confirmed prediction.

\textbf{The spirit of GTE predictions:} Every new measurement by Belle~II, DUNE,
CMB-S4, Hyper-Kamiokande, the Einstein Telescope, or Euclid is simultaneously a
test of GTE.  The framework does not hide in parameter space; it makes a
commitment at each observable.  The complete falsification register is presented
in the \emph{Falsifiability} tutorial~\cite{SpivackTutorialFalsifiability}.

%══════════════════════════════════════════════════════════════════════════════
\section{Summary: The Complete New-Physics Picture}
\label{sec:summary}
%══════════════════════════════════════════════════════════════════════════════

\begin{table}[htbp]
\centering
\caption{New-physics predictions and structural absences of GTE.  ``Absence'' means
structurally forbidden by the certificate; ``Prediction'' means explicitly derived.}
\label{tab:newphys}
\setlength{\tabcolsep}{5pt}
\begin{tabular}{lp{5.0cm}p{4.2cm}}
\toprule
Item & GTE position & Key experiment \\
\midrule
Supersymmetric particles & Absent (no winding sector) & LHC (confirmed null) \\
Axion & Absent ($\theta_{\rm QCD}=0$ from $F_{21}$) & ADMX, CASPEr (confirmed null) \\
Dark photon & Absent (no U(1) in dark sector) & NA64, Belle II (confirmed null) \\
Large extra dimensions & Absent ($3{+}1$D forced) & LHC (confirmed null) \\
\midrule
GTE-P7 at 211.9 MeV & Predicted (\CatAD) & \textbf{Belle II} \\
$\chi_1$ dark matter at 0.54 MeV & Predicted (\CatAD) & Sub-MeV DM experiments \\
$\chi_3$ at 3.60 GeV & Predicted (\CatAD) & LHC mono-jet, Higgs invisible \\
Dim-4 proton decay: none & Predicted (\CatAL) & Hyper-K, DUNE \\
$r = 0$ primordial tensors & Predicted (\CatA) & CMB-S4 ($\sim$2028) \\
$\mathbb{Z}_7$ GW background & Predicted (\CatAD) & Einstein Telescope \\
$\Omega_{\rm DM} h^2 = 0.11994$ & Predicted (\CatAD) & Planck ($-0.15\%$) \\
$\Sigma m_\nu = 59.4$~meV & Predicted (\CatA) & CMB-S4, Euclid \\
\bottomrule
\end{tabular}
\end{table}

The new-physics picture from GTE is internally consistent and makes a clear
experimental commitment: a dark sector with specific masses (the mirror branch),
proton stability at dimension-4, no axion, no SUSY, and a specific GW background.
Each prediction carries a named falsification criterion.

\section*{Further Reading}
\begin{tcolorbox}[colback=orange!8,colframe=orange!50!black,
  title=\textbf{Primary Sources}]
The results in this tutorial are derived in the following papers,
all available open-access on Zenodo and at
\href{https://novaspivack.com/research}{novaspivack.com/research}:
\begin{itemize}[itemsep=2pt]
  \item \textbf{P48 --- The Complete GTE Framework} (main synthesis monograph):\\
        \href{https://doi.org/10.5281/zenodo.20560550}{doi.org/10.5281/zenodo.20560550}
  \item \textbf{P47 --- Cosmological Predictions of the GTE/$\Phi_{\rm MDL}$ Framework}
        (cosmological predictions including $n_s$, $r$, and the GW background):\\
        \href{https://doi.org/10.5281/zenodo.20465803}{doi.org/10.5281/zenodo.20465803}
  \item \textbf{P08 --- From UGP to GTE: Prime-Locked Universes}
        (mirror-branch dark sector derivation):\\
        \href{https://doi.org/10.5281/zenodo.20171002}{doi.org/10.5281/zenodo.20171002}
  \item \textbf{P28 --- Computational Universality and the Standard Model}
        (no-CA-replica theorem, dimensional derivation):\\
        \href{https://doi.org/10.5281/zenodo.20259513}{doi.org/10.5281/zenodo.20259513}
\end{itemize}
Lean~4 machine proofs:
\href{https://github.com/novaspivack/ugp-lean}{github.com/novaspivack/ugp-lean}
\end{tcolorbox}

\bibliography{../../papers/bib/Spivack_Papers_Bibliography}
\end{document}
